\subsection{Irreducible spaces}

DEFINITION 1. A topological space X is called irreducible if every finite intersection of non-empty open sets of X is non-empty.

By considering the empty family of open sets of X it is seen that an irreducible space is non-empty; for a topological space X to be irreducible, it is necessary and sufficient that it be non-empty and that the intersection of two non-empty open sets of X be always non-empty (or, what amounts to the same thing, that the union of two closed sets distinct from X be always distinct from X).

PROPOSITION 1. Let X be a non-empty topological space. The following conditions are equivalent:

\begin{enumerate}
\def\labelenumi{(\alph{enumi})}
\item
  \(\mathbf{X}\) is irreducible;
\item
  every non-empty open set of \(\mathbf{X}\) is dense in \(\mathbf{X}\) ;
\item
  every open set of \(\mathbf{X}\) is connected.
\end{enumerate}

By definition, a set which is dense in X is a set which meets every non-empty open set, hence (a) and (b) are equivalent. It is immediate that (c) implies (a), for if \(U_{1}\) and \(U_{2}\) are disjoint non-empty open sets, \(U_{1} \cup U_{2}\) is a disconnected open set. Finally let us show that (a) implies (c): if U is a disconnected open set, it is the union of two disjoint non-empty sets \(U'\) , \(U''\) which are open in U and hence also open in X, which implies that X is not irreducible.

A Hausdorff space is irreducible only if it consists of a single point.

A subset E of a topological space X is called an irreducible set if the subspace E of X is irreducible. For this to be so, it is necessary and sufficient that, for every pair of sets U, V which are open in X and meet E, \(U \cap V\) also meet E, or (what amounts to the same thing) that, for every pair of sets F, G which are closed in X and satisfy \(E \subset F \cup G\) , \(E \subset F\) or \(E \subset G\) . By induction on n, we deduce that, if \((\mathbf{F}_{i})_{1 \leq i \leq n}\) is a finite family of closed sets of X such that \(E \subset \bigcup_{i=1}^{n} F_{i}\) , there exists an index i such that \(E \subset F_{i}\) .

PROPOSITION 2. In a topological space X, for a set E to be irreducible, it is necessary and sufficient that its closure \(\overline{E}\) be so.

For an open set of X to meet E, it is necessary and sufficient that it meet \(\overline{E}\) ; then the proposition follows from the above remarks.

PROPOSITION 3. (i) If X is an irreducible space, every non-empty open subset of X is irreducible.

\begin{enumerate}
\def\labelenumi{(\roman{enumi})}
\setcounter{enumi}{1}
\item
  Let \((\mathbf{U}_{\alpha})_{\alpha \in \mathbf{A}}\) be a non-empty covering of a topological space \(\mathbf{X}\) consisting of open sets such that \(\mathbf{U}_{\alpha} \cap \mathbf{U}_{\beta} \neq \varnothing\) for every pair of indices \((\alpha, \beta)\) . If the sets \(\mathbf{U}_{\alpha}\) are irreducible, the space \(\mathbf{X}\) is irreducible.
\end{enumerate}

(i) If \(\mathbf{X}\) is irreducible, \(\mathbf{U} \subset \mathbf{X}\) non-empty and open in \(\mathbf{X}\) and \(\mathbf{V} \subset \mathbf{U}\) non-empty and open in \(\mathbf{U}\) , \(\mathbf{V}\) is also open in \(\mathbf{X}\) , therefore dense in \(\mathbf{X}\) and a fortiori dense in \(\mathbf{U}\) . Then \(\mathbf{U}\) is irreducible (Proposition 1).

(ii) Let us show that, for every non-empty open set V in X, \(V \cap U_{\alpha} \neq \varnothing\) for all \(\alpha \in A\) : it follows that \(V \cap U_{\alpha}\) is dense in \(U_{\alpha}\) by hypothesis, hence that V is dense in X and this proves that X is irreducible (Proposition 1). Now there exists at least one index \(\gamma\) such that \(V \cap U_{\gamma} \neq \varnothing\) ; as \(U_{\alpha} \cap U_{\gamma} \neq \varnothing\) for all \(\alpha\) and \(V \cap U_{\gamma}\) is dense in \(U_{\gamma}\) , \(U_{\alpha} \cap U_{\gamma} \cap V \neq \varnothing\) and a fortiori \(U_{\alpha} \cap V \neq \varnothing\) , which completes the proof of (ii).

PROPOSITION 4. Let X and Y be two topological spaces and f a continuous mapping from X to Y. For every irreducible subset E of X, \(f(E)\) is an irreducible subset of Y.

If U, V are two open sets of Y which meet \(f(\mathrm{E})\) , \(\overline{f}^{-1}(\mathrm{U})\) and \(\overline{f}^{-1}(\mathrm{V})\) are open sets of X which meet E. Consequently, \(\overline{f}^{-1}(\mathrm{U}) \cap \overline{f}^{-1}(\mathrm{V}) = \overline{f}^{-1}(\mathrm{U} \cap \mathrm{V})\) meets E, which implies that \(U \cap V\) meets \(f(\mathrm{E})\) and proves the proposition.

DEFINITION 2. Every maximal irreducible subset of a topological space X is called an irreducible component of X.

It follows from Proposition 2 that every irreducible component of \(\mathbf{X}\) is closed in \(\mathbf{X}\) .

PROPOSITION 5. Let X be a topological space. Every irreducible subset of X is contained in an irreducible component of X and X is the union of its irreducible components.

To prove the first assertion, it is sufficient, by virtue of Zorn's Lemma, to prove that the set \(\Im\) of irreducible subsets of \(X\) is inductive. Let \(\mathfrak{G}\) be a subset of \(\Im\) totally ordered by inclusion; we show that the union \(E\) of the sets \(F \in \mathfrak{G}\) is irreducible. Let \(U\) , \(V\) be two open sets of \(X\) which meet \(E\) ; as \(\mathfrak{G}\) is totally ordered, there exists a set \(F \in \mathfrak{G}\) meeting U and V; as F is irreducible, \(U \cap V\) meets F and hence also E, which proves that E is irreducible and hence that \(\mathfrak{G}\) is inductive. The second assertion follows from the first, for every subset of X consisting of a single point is irreducible.

COROLLARY. Every connected component of a topological space X is a union of irreducible components of X.

Every irreducible subspace of X is connected by Proposition 1 and hence is contained in a connected component of X.

Note that two distinct irreducible components of X may have points in common (Exercise 11).

PROPOSITION 6. Let X be a topological space and \((\mathrm{P}_{i})_{1\leq i\leq n}\) a finite covering of X consisting of irreducible closed sets. Then the irreducible components of X are the maximal elements (with respect to inclusion) of the set of \(P_{i}\) .

We may restrict ourselves to the case where no two \(P_{i}\) are comparable. If E is an irreducible subset of X, then \(E \subset \bigcup_{i=1}^{n} P_{i}\) and hence E is contained in one of the closed sets \(P_{i}\) ; this proves that the \(P_{i}\) are the only maximal irreducible subsets of X.

COROLLARY. Let X be a topological space and E a subspace with only a finite number of distinct irreducible components \(Q_{i}\) ( \(1 \leqslant i \leqslant n\) ); then the irreducible components of the closure \(\overline{E}\) in X are the closures \(\overline{Q}_{i}\) of the \(Q_{i}\) ( \(1 \leqslant i \leqslant n\) ) and \(\overline{Q}_{i} \neq \overline{Q}_{j}\) for \(i \neq j\) .

\(\overline{\mathbf{E}}\) is the union of the \(\overline{\mathbf{Q}}_i\) , which are irreducible (Proposition 2); as \(\mathbf{Q}_i\) is closed in \(\mathbf{E}\) , \(\overline{\mathbf{Q}}_i \cap \mathbf{E} = \mathbf{Q}_i\) ; as \(\mathbf{Q}_i \not\subset \mathbf{Q}_j\) for \(i \neq j\) , \(\overline{\mathbf{Q}}_i \not\subset \overline{\mathbf{Q}}_j\) , whence the corollary by virtue of Proposition 6.

Remark. Suppose that X has only a finite number of distinct irreducible components \(X_{i}\) ( \(1 \leqslant i \leqslant n\) ); then \(U_{i} = \complement\left(\bigcup_{j \neq i} X_{j}\right)\) is open in X and dense in \(X_{i}\) since \(X_{i} \not\subset \bigcup_{j \neq i} X_{j}\) ; the \(U_{i}\) ( \(1 \leqslant i \leqslant n\) ) are therefore non-empty open sets of X which are irreducible (Proposition 2), pairwise disjoint and with their union dense in X.

PROPOSITION 7. Let U be an open subset of a topological space X. The mapping \(V \mapsto \overline{V}\) (closure in X) is a bijection of the set of irreducible subsets of U which are closed in U onto the set of irreducible subsets of X which are closed in X and meet U; the inverse bijection is \(Z \mapsto Z \cap U\). In particular, this bijection maps the set of irreducible components of U onto the set of irreducible components of X which meet U.

If V is closed in U and irreducible, \(\overline{V}\) is irreducible (Proposition 2) and \(V = \overline{V} \cap U\) . Conversely, if Z is irreducible and closed in X and meets U, \(Z \cap U\) is a non-empty open subset of Z, hence irreducible (Proposition 3) and dense in Z and, as Z is closed, \(Z = \overline{Z \cap U}\) . This proves the proposition.

